\section{实验与验证}

\subsection{实验设计原则}

每个 DSPy experiment 应明确 `hypothesis`、`metric`、`baseline`。推荐 three-phase pipeline：1) small-scale data → 2) evidence capture → 3) production dry-run。跨团队分享 evidence spreadsheet 让别的团队复现实验。

\subsection{Simulation to Production}

Experimental validation 不能止步于 offline benchmark，还需要 `sim-to-prod` steps。每次 optimizer change 都要在 sandbox environment 运行 1000+ queries，并记录 `metric drift` 与 `cost per query`，确定 rollout 可控。

\begin{knowledgebox}{Simulation Pipeline}
Simulation pipeline 包括 `toy dataset` run、`evidence replay`、`sanity check`，每一步都有 defined owner。若 evidence replay 结果差异 >5\%，必须重新 tune 或 abort。
\end{knowledgebox}

\subsection{本章小结}

实验与验证段落帮助团队把 DSPy 调整变得可复现，并在 production rollout 前建立信心。

